\section{Arithmetic, names, and certificates}\label{sec:arithmetic}

\subsection{The base theory}

We use the usual definitional presentation of intuitionistic primitive
recursive arithmetic. Its only object sort is $\N$. The language contains
$0$, successor, and a function symbol for each fixed primitive recursive
function, with the defining equations for composition and primitive
recursion. Equality and the ordinary intuitionistic first-order connectives
and quantifiers are available. The nonlogical axioms include successor
injectivity, $Sx\ne0$, the defining equations, and induction
\[
 A(0,\vec a)\ \wedge\
 \forall n\bigl(A(n,\vec a)\to A(n+1,\vec a)\bigr)
 \ \longrightarrow\ \forall n\,A(n,\vec a)
\]
for quantifier-free $A$. A presentation with fewer primitive symbols can
instead use its standard definitional extensions.

We use no excluded middle for arbitrary formulas, choice principle,
comprehension principle, Markov principle, or induction on arbitrary
quantified analytic formulas. Rational comparison is decidable, and the
corresponding finite case distinctions are permitted intuitionistically.

\subsection{Concrete number coding}

One possible signed-integer code sends $2a$ to $a$ and $2a+1$ to $-(a+1)$.
A rational is a pair $(z,d)$ representing the integer coded by $z$ divided
by $d+1$. Fractions need not be reduced. Pairing can use
\[
 \langle a,b\rangle=\frac{(a+b)(a+b+1)}2+b.
\]
Its projections are obtained by bounded search; the displayed division is
exact. Finite vectors are nested pairs together with their length.

Addition, multiplication, sign reversal, absolute value, and rational order
are primitive recursive. To compare $a/b$ with $c/d$, with $b,d>0$, compare
$ad$ and $cb$ as signed integers. All denominators are positive by
construction. Division by a known positive integer, powers, bounded
products, finite sums, finite maxima, and grid-index decoding are therefore
primitive recursive too.

We write $\epsilon_p=2^{-p}$ as a rational abbreviation. For an integer
$A\ge1$, write $\ell(A)=\lceil\log_2 A\rceil$. This is the bounded search
for the least $j\le A$ with $A\le2^j$. We put $\ell(0)=0$ when needed.
For a modulus $\mu$ use the monotone majorant
\[
 \bar\mu(t)=\max_{s\le t}\mu(s).
\]
This is a bounded recursion, not an unbounded search.

\subsection{A schema, not a universal evaluator}\label{sec:schema}

An evaluation function such as $F_i(x,p)$ is fixed when an instance of the
theorem is formed. Numerical parameters in one such fixed function may vary
uniformly. The theorem may also use one fixed primitive recursive family
$F(a,i,x,p)$ with number parameter $a$. It does not quantify internally over
all functions $\N\to\N$.

In particular, there is no primitive recursive function $U(e,x)$ which
correctly evaluates every primitive recursive function indexed by $e$.
If there were, $x\mapsto U(x,x)+1$ would itself have an index $d$, giving
$U(d,d)=U(d,d)+1$. Nothing in this report assumes such a universal evaluator.
Compiling a \emph{fixed} syntactic program into terms is different from
claiming a total primitive recursive evaluation function for every program.

\subsection{Regular names and equality}

\begin{definition}[Regular rational name]
A regular name is a fixed primitive recursive function $a:\N\to\Q$ with
the arithmetic property
\begin{equation}\label{eq:regular}
 \forall p,q,\qquad \abs{a(p)-a(q)}\le\epsilon_p+\epsilon_q.
\end{equation}
Two regular names are equal, written $a\EqName b$, if
\begin{equation}\label{eq:name-equality}
 \forall p,q,\qquad \abs{a(p)-b(q)}\le\epsilon_p+\epsilon_q.
\end{equation}
These are first-order arithmetic formulas after the fixed programs and
rational coding have been expanded.
\end{definition}

For regular names, the order abbreviation is
\[
 a\le b\quad:\Longleftrightarrow\quad
 \forall p,q\ \bigl(a(p)\le b(q)+\epsilon_p+\epsilon_q\bigr).
\]
A rational constant is represented by its constant name. Absolute value
and finite maximum use their rational evaluators, which are
$1$-Lipschitz for the sup norm. Sums and products use the guarded evaluators
in \cref{sec:calculus}. These conventions specify the named norms and
radius comparisons used in the jet arguments. They do not provide a
decision procedure for named-real order.

For explanatory purposes a regular name may be pictured as converging to a
real number with error $\epsilon_p$ at precision $p$. The proof below never
needs to construct a completed real-number object inside $\ipra$.

\begin{lemma}[Eliminating vanishing rational slack]\label{lem:slack}
For rationals $x,a$ and an integer $C\ge0$, if
$x\le a+C2^{-q}$ for every $q$, then $x\le a$. This argument is available in
$\ipra$.
\end{lemma}
\begin{proof}
Decide the rational comparison $x\le a$. In the other case write
$x-a=u/v$ with positive integers $u,v$. Set $q=C+v+1$.
The elementary induction $2^n\ge n+1$ gives
$2^q=2\cdot2^C2^v\ge2(C+1)(v+1)>Cv$.
Consequently $C2^{-q}<1/v\le x-a$, contrary to the hypothesis at this
particular $q$. All comparisons are rational comparisons and $q$ is a
primitive recursive witness.
\end{proof}

Reflexivity and symmetry of \eqref{eq:name-equality} are immediate.
For transitivity, insert a sufficiently precise sample of the middle name,
obtain the desired inequality with slack $2^{1-q}$, and apply
\cref{lem:slack}. Thus equality is an equivalence relation, without deciding
equality or order of represented real numbers.

A convenient stronger witness for equality is a primitive recursive
function $E$ such that
\begin{equation}\label{eq:tail-equality}
 \forall r\ \forall a,b\ge E(r),\quad
 \abs{A(a)-B(b)}\le2^{-r}.
\end{equation}
Together with regularity this implies \eqref{eq:name-equality}: compare
$A(p),B(q)$ to samples of common index
$m\ge\max\{E(r),r\}$ and remove the residual $3\epsilon_r$ by
\cref{lem:slack}. The concrete Stokes names satisfy stronger constants
that also give \eqref{eq:name-equality} directly.

\subsection{The analytic input on a rational cube}\label{sec:certificates}

Fix $k\ge1$. For every $i$, the programs $F_i,G_i$ return rational values
at rational points of $Q^k$. Supply primitive recursive moduli
$\mu_F,\mu_G,\tau$. They may already be monotone; otherwise use their
bounded-max majorants.

\paragraph{Coherence (C).}
For $H=F_i$ and $H=G_i$,
\begin{equation}\label{eq:coherence}
 \abs{H(x,p)-H(x,q)}\le\epsilon_p+\epsilon_q.
\end{equation}

\paragraph{Uniform continuity (U).}
For $H=F_i$ use $\mu_H=\mu_F$; for $H=G_i$ use $\mu_H=\mu_G$.
For all rational $x,y\in Q^k$ and all $t,p,q$,
\begin{equation}\label{eq:uc}
 \norm{x-y}\le2^{-\mu_H(t)}
 \quad\Longrightarrow\quad
 \abs{H(x,p)-H(y,q)}
 \le\epsilon_t+\epsilon_p+\epsilon_q.
\end{equation}
The norm is the maximum of the finitely many absolute coordinate
differences. A common modulus for a finite family is obtained by a bounded
maximum.

\paragraph{Uniform coordinate remainder (D).}
For rational $h>0$ with $x,x+he_i\in Q^k$,
\begin{equation}\label{eq:derivative-cert}
 h\le2^{-\tau(t)}
 \quad\Longrightarrow\quad
 \abs{F_i(x+he_i,p)-F_i(x,p)-hG_i(x,p)}
 \le h\epsilon_t+(2+h)\epsilon_p.
\end{equation}
Here $e_i$ is the $i$th coordinate vector.

The quantified variables in these conditions are number codes for rational
points and precisions. Invalid codes or points outside the box are excluded
by decidable guards. We denote their conjunction by
$\mathsf H(F,G,\mu_F,\mu_G,\tau)$.

\begin{remark}[The remainder assumption is local]
Condition \eqref{eq:derivative-cert} is a uniform first-order approximation
property. It contains no integral or boundary identity. If exact named
values satisfy
$|f_i(x+he_i)-f_i(x)-hg_i(x)|\le h\epsilon_t$ and all three evaluated
quantities have error $\epsilon_p$, the two coefficient errors and the
derivative error give exactly $(2+h)\epsilon_p$.
Only the diagonal derivatives $g_i=\partial_i f_i$ are needed for the cube
theorem. Full jets are used later to produce these data from curved
parameterizations.
\end{remark}

\subsection{Finite arithmetic versus universal certificates}

The theorem is conditional on $\mathsf H$. A proof for a particular analytic
object must supply proofs of these conditions; a program returning decimal
values is insufficient. For polynomial objects those proofs are obtained
by finite binomial identities in \cref{sec:polynomial-jets}. Other certified
analytic libraries can supply further examples.

For any one requested output precision the proof of the finite estimate uses
only finitely many instances of (D), and the refinement estimate uses only
finitely many instances of (U). If desired those instances can be collected
into a primitive recursive Boolean fold, a finite certificate for that
request. Universal correctness of the evaluator or of all precisions is
not decided by this finite checker.
